\section{作品测试与分析}

本章围绕消息恢复、密钥控制、密文图像的视觉质量、跨设备一致性及生成效率，介绍本作品的测试方案、过程与结果。

\subsection{测试方案}

测试以一条消息的生成与恢复为基本流程：发送方选取初始图像和预共享图像密钥，将目标消息编码为神经元分组状态，经迭代优化得到密文图像；接收方使用相同密钥、DNN 和编码配置恢复消息。首先以正确密钥检查功能正确性，再更换随机密钥、其他样本的密钥或不提供密钥，观察恢复结果。为检验生成的随机性，在固定消息与密钥的条件下，分别改变初始图像及优化随机状态，或只改变优化随机状态，比较多次生成的图像。

测试还从四个方面考察作品性能。其一，比较不同分组大小下的生成成功率、迭代次数和用时，判断分组编码对不响应神经元的缓解作用。其二，以峰值信噪比（PSNR）、结构相似性（SSIM）、像素差异及直接泄漏代理指标评价图像质量，并比较启用和关闭间接泄漏损失时的神经元状态偏差。其三，测试截获样本分析、生成时间分析、跨设备恢复及图像失真后的恢复。其四，在相同消息长度和图像尺寸下比较生成、恢复及模型准备成本。消息恢复采用比特错误率（BER）和完整消息恢复率（EMR）：BER 越低越好，EMR 表示零比特错误消息所占比例。图像失真测试只报告 BER，不能据此推断每条消息均可完整恢复。

\subsection{测试环境搭建}

实验使用 ImageNet 数据集和实拍图像，图像尺寸包括 $224\times224$ 与 $896\times896$。实验在 Ubuntu 系统上以 Python 3.10 和 PyTorch 运行，使用 RTX 4090 GPU。默认选取 ShuffleNet V2 的 softmax 前分类分数作为编码所用的 DNN 响应，利用随机化快速梯度符号方法（R-FGSM）生成密文图像。直接与间接泄漏损失的权重分别设为 $\alpha=1000$ 和 $\beta=100$，分组大小 $g=2$，编码阈值为 $0.5$，学习率为 $10^{-4}$。除专门比较分组、学习率、模型或设备的实验外，其余测试沿用上述配置。

生成端和恢复端需要采用相同的模型、神经元选择、分组方式、阈值和数值处理规则，并预先秘密共享图像密钥。跨设备测试另外比较了启用与关闭截断处理 $T_{\mathrm{trc}}$ 的结果，以观察计算精度差异对消息恢复的影响。

\subsection{测试设备}

默认实验平台采用 RTX 4090 GPU。跨设备兼容性实验还使用 RTX 3070 GPU、Intel Core i7-13700K CPU 和 AMD Ryzen 9 5900HX CPU，并组合不同的生成端与恢复端。该实验用于检验同一消息和初始图像在不同计算设备之间传递时，数值精度变化是否造成解码错误。设备组合及结果见本章第 3.4 节。

\subsection{测试功能}

\textbf{消息生成与密钥控制。}本作品使用 100 组由不同图像密钥生成的消息与密文图像测试恢复功能。使用正确密钥时，BER 为 $0\%$、EMR 为 $100\%$；随机密钥、其他样本的错误密钥和无密钥时，BER 分别为 $49.9\%$、$49.7\%$ 和 $50.1\%$，EMR 均为 $0\%$。因此，在这组测试中，只有正确密钥实现了完整恢复；错误或缺失密钥下的逐比特结果接近随机猜测。

\textbf{概率生成。}固定消息和图像密钥，分别独立采样初始图像及 R-FGSM 随机状态，完整生成 30 次；再固定初始图像，仅改变优化随机状态，生成 30 次。两组测试均未出现相同的固定精度输出。图~\ref{fig:probabilistic-cipher-images} 展示了第二组实验的代表性密文图像及其与原始图像的结构相似度，平均 SSIM 约为 $97.5\%$。这说明在所测样本中，内部随机优化也能改变输出，同时保持较高的视觉相似度；有限次数的未碰撞观察不能给出任意运行次数下的碰撞概率。

\begin{figure}[!htbp]
\centering
\IfFileExists{figures/cipher-images-original-panel.png}{%
  \def\cipherpanelpath{figures/cipher-images-original-panel.png}%
}{%
  \def\cipherpanelpath{C:/Users/Li Zhe/Desktop/SSIR_report-main/figures/cipher-images-original-panel.png}%
}
\begin{tikzpicture}
\node[anchor=south west,inner sep=0] (panel) at (0,0)
  {\includegraphics[width=0.86\textwidth,trim=0 65 0 0,clip]{\cipherpanelpath}};
\begin{scope}[x={(panel.south east)},y={(panel.north west)}]
  \fill[white] (0,0) rectangle (1,0.09);
  \fill[white] (0,0.49) rectangle (1,0.55);
  \node[font=\scriptsize] at (0.164,0.52) {相似度：97.46\%};
  \node[font=\scriptsize] at (0.486,0.52) {相似度：97.45\%};
  \node[font=\scriptsize] at (0.807,0.52) {相似度：97.48\%};
  \node[font=\scriptsize] at (0.164,0.045) {相似度：97.50\%};
  \node[font=\scriptsize] at (0.486,0.045) {相似度：97.48\%};
  \node[font=\scriptsize] at (0.807,0.045) {原始图像};
\end{scope}
\end{tikzpicture}
\caption{固定初始图像时生成的密文图像与原始图像}
\label{fig:probabilistic-cipher-images}
\end{figure}

图~\ref{fig:probabilistic-cipher-images} 中，五幅密文图像使用相同的初始图像、消息和图像密钥，仅改变优化过程的随机状态；右下角为原始图像。五幅密文图像的 SSIM 分别为 $97.46\%$、$97.45\%$、$97.48\%$、$97.50\%$ 和 $97.48\%$，数值集中在 $97.5\%$ 附近。从图中整体观察，物体轮廓、位置和主要颜色均与原图接近，说明这些生成结果在保持图像内容方面表现稳定；与此同时，五次输出并非同一幅固定精度图像。图示反映的是其中一次固定初始图像的代表性结果，多次生成是否重复仍应结合前述两组各 30 次的完整测试判断。SSIM 衡量视觉结构相似性，不能单独说明消息恢复是否正确，也不能作为密文安全性的证明。

\textbf{分组编码。}在 EfficientNet-B0 的 1280 维预分类特征上，以 150 位消息测试 $g\in\{1,2,4,6\}$。每种设置进行了 2000 次配对试验，最大迭代次数为 3000。直接编码（$g=1$）的成功率为 $79\%$，$g=2,4,6$ 均为 $100\%$。其中 $g=2$ 的成功样本平均迭代 $47.5$ 次、生成耗时 $317.8$~ms；$g=4$ 和 $g=6$ 分别为 $50.3$ 次、$336.5$~ms 及 $59.1$ 次、$389.7$~ms。在这一配置下，$g=2$ 同时保持较高表示容量和较低生成成本。

\textbf{图像质量与泄漏抑制。}以 $L_1$、$L_2$、SSIM 和 PSNR 衡量密文图像与原图的差异。本作品将 $(1-\mathrm{SSIM})\times\mathrm{MSE}$ 定义为直接视觉泄漏的代理指标；它衡量可见差异，不能等同于密码学区分优势。各方法的直接泄漏指标与图像质量见表~\ref{tab:direct-leakage}。学习率分别为 $10^{-3}$、$10^{-4}$、$10^{-5}$ 时，该指标为 $39.34$、$0.41$、$0.08$，对应 PSNR 为 $26.62$、$38.14$、$41.93$~dB，SSIM 为 $72.24\%$、$95.89\%$、$97.99\%$。

\begin{table}[!htbp]
\centering
\caption{直接信息泄漏对比}
\label{tab:direct-leakage}
\begin{tabular}{lccc}
\hline
方法 & PSNR & SSIM & 直接泄漏指标 \\
\hline
FNNS~\supercite{19-kishore-et-al-2021} & 33.01~dB & 91.56\% & 2.74 \\
SFNNS~\supercite{8-luo-et-al-2023} & 40.88~dB & 97.92\% & 0.11 \\
SteganoGAN-D~\supercite{4-zhang-et-al-2019-steganogan} & 38.40~dB & 91.33\% & 0.81 \\
SteganoGAN-B~\supercite{4-zhang-et-al-2019-steganogan} & 24.21~dB & 81.67\% & 45.17 \\
SSIR（$\mathrm{lr}=10^{-3}$） & 26.62~dB & 72.24\% & 39.34 \\
SSIR（$\mathrm{lr}=10^{-4}$） & 38.14~dB & 95.89\% & 0.41 \\
SSIR（$\mathrm{lr}=10^{-5}$） & 41.93~dB & 97.99\% & 0.08 \\
\hline
\end{tabular}
\end{table}

表~\ref{tab:direct-leakage} 将图像质量与直接泄漏代理指标并列，便于观察同一设置下视觉差异的变化。对本作品而言，学习率从 $10^{-3}$ 降至 $10^{-4}$ 时，PSNR 从 $26.62$~dB 升至 $38.14$~dB，SSIM 从 $72.24\%$ 升至 $95.89\%$，直接泄漏指标由 $39.34$ 降至 $0.41$；继续降至 $10^{-5}$ 后，三项结果分别变为 $41.93$~dB、$97.99\%$ 和 $0.08$。因此，表中差异主要来自学习率的选择，不能只报出其中一行作为所有配置的图像质量。按本章采用的定义，直接泄漏指标同时包含结构差异与均方误差，数值越小表示该代理指标下的可见差异越低。

与表中其他方法相比，本作品在 $\mathrm{lr}=10^{-3}$ 时的直接泄漏指标高于 FNNS、SFNNS 和 SteganoGAN-D；在 $\mathrm{lr}=10^{-5}$ 时则低于这些列出的数值。这里比较的是表中的图像质量和代理指标，并未涵盖各方法的消息容量、恢复条件或模型准备成本，也不意味着数值较小就具有更强的密码学安全性。学习率降低带来的视觉质量改善还需与生成迭代及耗时一并权衡。

图~\ref{fig:psnr-ssim-ccdf} 展示了不同学习率下密文图像 PSNR 与 SSIM 的互补累积分布。

\begin{figure}[!htbp]
\centering
\IfFileExists{figures/psnr-ssim-ccdf-zh.png}{%
  \includegraphics[width=\textwidth]{figures/psnr-ssim-ccdf-zh.png}%
}{%
  \includegraphics[width=\textwidth]{C:/Users/Li Zhe/Desktop/SSIR_report-main/figures/psnr-ssim-ccdf-zh.png}%
}
\caption{不同学习率下密文图像 PSNR 与 SSIM 的互补累积分布}
\label{fig:psnr-ssim-ccdf}
\end{figure}

图~\ref{fig:psnr-ssim-ccdf} 的纵轴表示达到或超过横轴阈值的密文图像比例，而不是单幅图像的 PSNR 或 SSIM。左图中，学习率为 $10^{-3}$ 的 PSNR 曲线最先下降，$10^{-4}$ 与 $10^{-5}$ 的曲线依次向右移动；右图的 SSIM 曲线也呈现相同的整体次序。这意味着降低学习率后，获得较高 PSNR 和 SSIM 的图像在所测样本中占比提高。曲线还展示了各配置内部的分布：同一学习率下，不同样本并非都具有表~\ref{tab:direct-leakage} 所列的平均质量，不能用一个平均值代替全部生成结果。

两幅图分别从像素误差和结构相似性观察图像质量，方向一致，但各自反映的差异并不完全相同。结合表~\ref{tab:direct-leakage}，较小学习率在这组测试中同时对应更高的图像质量和更低的直接泄漏代理指标。图中的曲线只能说明这些评价量在已测图像上的分布，无法推断消息是否一定可恢复，也不能据此给出未测试图像或其他参数下的保证。

图~\ref{fig:indirect-leakage} 比较了 ResNet50、DenseNet121、EfficientNet V2 和 Inception V3 在启用与关闭间接泄漏损失时的神经元状态偏差。四个模型的平均降幅为 $56.02\%$。

\begin{figure}[!htbp]
\centering
\IfFileExists{figures/indirect-leakage-zh.png}{%
  \includegraphics[width=\textwidth]{figures/indirect-leakage-zh.png}%
}{%
  \includegraphics[width=\textwidth]{C:/Users/Li Zhe/Desktop/SSIR_report-main/figures/indirect-leakage-zh.png}%
}
\caption{四种网络的间接信息泄漏对比}
\label{fig:indirect-leakage}
\end{figure}

图~\ref{fig:indirect-leakage} 按网络分别列出若干层及其总体状态偏差，蓝色表示未使用间接泄漏损失，青色表示使用该损失。启用损失后，四个网络所标出的平均状态偏差都下降：ResNet50 从 $0.179$ 降至 $0.069$，DenseNet121 从 $0.150$ 降至 $0.058$，EfficientNet V2 从 $0.523$ 降至 $0.252$，Inception V3 从 $0.117$ 降至 $0.047$。其中 EfficientNet V2 的原始偏差及下降后的偏差都高于其余三个网络，说明不同网络对这项约束的响应幅度并不相同。

从各层柱形看，较大的偏差集中在少数层或分类端，例如 ResNet50 的 Fc、DenseNet121 的 Class 以及 EfficientNet V2 的 L2 和 L6；加入损失后，这些位置的柱形均降低。因而，平均值的改善不是由所有层以相同幅度变化形成的，逐层观察有助于识别仍需关注的位置。这里的状态偏差用于评估间接泄漏抑制效果，不等同于从密文图像恢复消息的成功率；图示也只覆盖所列四种网络及其选取的层。

\textbf{攻击、设备与传输扰动。}截获消息--密文图像配对样本后训练恢复模型的实验中，随着样本数量增加，BER 仍约为 $50\%$，未呈现稳定下降趋势。对高斯噪声图像密钥改变标准差并记录生成时间的实验，未观察到生成时间与密钥统计量之间的一致关系。这两项结论只适用于所测攻击方式和参数范围。跨设备测试中，启用截断处理后，RTX 4090、RTX 3070、i7-13700K 和 Ryzen 9 5900HX 的已测生成/恢复组合均达到 $0\%$ BER；未启用时，部分组合出现 $0.67\%$ 至 $2.6\%$ 的 BER。对密文图像施加高斯噪声、椒盐噪声、散斑噪声或 JPEG 压缩后，所测设置的 BER 为 $0.03\%$ 至 $0.82\%$，表明其具有一定的比特级失真容忍度，但不能保证有损传输后的完整消息恢复。

\subsection{结果分析}

测试结果说明，SSIR 在上述配置下实现了由图像密钥控制的消息恢复。正确密钥测试达到零比特错误，而错误密钥和无密钥测试接近随机猜测。分组编码将所测成功率从直接编码的 $79\%$ 提高到 $100\%$；继续增大分组未带来额外成功率收益，却降低了表示容量并增加生成时间。这支持选取 $g=2$ 作为默认配置。两组各 30 次的随机生成测试也显示，初始图像和优化过程均可影响密文图像的具体形式。

视觉质量与速度之间存在可见取舍。学习率由 $10^{-3}$ 降至 $10^{-5}$ 时，直接泄漏代理指标从 $39.34$ 降至 $0.08$，相对降幅为 $99.79\%$；PSNR 与 SSIM 同时提高，但更小的学习率通常需要更多迭代。150 位消息、$224\times224$ 图像的效率比较中，SSIR 的密文图像生成耗时约 $300$--$500$~ms，恢复约 $10$~ms，且无需专用模型准备；对照方法的嵌入耗时为 $50$--$120$~ms，但模型准备需数小时。因而，SSIR 的在线生成速度并非最快，其优势主要体现在复用现有 DNN、避免专用模型准备，并同时满足所测配置下的密钥控制与视觉质量要求。

测试也给出适用边界。跨设备零 BER 依赖已测硬件组合及截断处理；失真测试虽显示较低 BER，仍会出现比特错误。直接泄漏数值是视觉差异的代理指标，截获样本和时间分析覆盖的是特定攻击条件。形式化的 IND-CPA 与 IND-CCA1 结论还依赖输出分布条件和每密钥查询限制，不能仅凭上述实验数值推出无条件安全性。相关形式化分析见下一节，具体证明见附录。
\subsection{安全性分析}
实验测试表明了方案在模拟攻击下表现良好，但密码方案的安全性需要严谨的数学推导作为支撑。为此，我们从三个层次递进展开：
首先，分析了底层的DNN非线性响应机制，说明信息如何可靠嵌入并恢复。其次，刻画密文图像的概率生成过程与多语义特性，分析攻击者可能利用的随机性特征。
最后，定义标准的IND-CPA与IND-CCA1密码学安全游戏，并在限制查询次数的前提下，推导攻击者成功区分密文的概率上界。
本节中，ENC 表示算法 1 的完整生成流程（包括泄露损失与截断），DEC 表示消息恢复过程。
\subsubsection{用于密钥控制表示的DNN非线性}
本节主要阐明信息如何被可靠地嵌入DNA并准确提取。首先定义确定地映射函数$F_k$,然后划分融合图像区域，证明只要生成的图像落在特定区域内，就能通过阈值解码准确恢复消息，从而确立系统的恢复一致性。

\textbf{1) 输入依赖的DNN响应：}
SSIR使用选定神经元的分组响应来表示图像密钥下的消息。
对于不同的密文图像，添加相同的密钥会在这些连续响应中产生不同的变化。
因此，DNN的非线性响应可以使密钥诱导的变化依赖于当前图像。
附录A通过与操作域上的固定仿射映射进行比较，阐明了这种相互作用。

\textbf{2) DNN响应的离散使用：}
受密码操作神经实现[7]的启发，SSIR通过分组编码将连续DNN响应连接到离散消息状态。
令 \(\mathcal{M}\) 为支持的等长消息空间，具有双射编码 \(\mathrm{Et} : \mathcal{M} \to \{0,1\}^\ell\)，\(C_{\mathrm{fp}}\) 为有限固定精度密文图像域，\(\mathcal{K} = \mathrm{supp}(\mathrm{KeyGen})\) 为合法图像密钥域。
我们考虑支持的配置，其中对于每个 \(k \in \mathcal{K}\) 和 \(m \in \mathcal{M}\)，生成以正概率成功。
定义
\begin{equation}
\begin{array}{rcl}
F_k(c) &=& \mathrm{Ec}(z(c+k); g, t), \\
\mathrm{DEC}_k(c) &=& \mathrm{Et}^{-1}(F_k(c)).
\end{array}
\end{equation}

在约定的DNN、编码和数值设置下，双方评估相同的确定性映射 \(F_k\)。
对于 \(\mathcal{I}_{\mathrm{fp}} = C_{\mathrm{fp}} + \mathcal{K}\)，将与字符串 \(s\) 相关联的融合图像区域定义为
\begin{equation}
\Omega_s = \{ i \in \mathcal{I}_{\mathrm{fp}} : \mathrm{Ec}(z(i); g, t) = s \}.
\end{equation}
那么当且仅当 \(c + k \in \Omega_s\) 时，\(F_k(c) = s\)。
在生成过程中，连续分组响应提供优化信号，而离散条件 \(F_k(p_n) = \mathrm{Et}(m)\) 决定成功终止。
因此，所得的 \(c = p_n\) 满足 \(\mathrm{DEC}_k(c) = m\)。
因此，相同的DNN响应指导状态诱导，并在阈值处理后确定恢复的字符串。
当分组响应的变化超过配置的阈值时，它会影响所表示的比特。
附录B详细介绍了这种从连续到离散的构造及其恢复一致性。

通过刻画输入依赖的非线性相应并构建分组阈值编码，我们证明了确定性映射下的恢复一致性。这意味着，在约定的DNN、编码和数值精度下，只要通信双方使用正确的图像密钥，
系统在数学上即能准确恢复原始消息，为后续分析密文图像的安全边界提供了正确性基础。
\subsubsection{概率生成与多语义}

SSIR支持两种互补的表示形式：生成可以为固定消息和密钥产生不同的图像，而固定图像在兼容密钥下可以接受不同的消息解释。
我们引入初始图像和优化状态双重随机源，推导生成结果的概率分布并计算出碰撞概率上界。
同时，利用融合加法运算的可交换性质，构造出兼容密钥集合，证明单张图像可能对应多种合法解释。

\textbf{1) 概率生成：}
每次生成尝试使用两个独立采样的随机源：
\begin{equation}
r \leftarrow D_R, \quad \xi \leftarrow D_\Xi, \quad c \leftarrow \mathrm{ENC}_k(m; r, \xi). \tag{13}
\end{equation}
初始图像 \(r\) 确定 \(p_0 = r/255\) 以及泄漏损失所使用的参考。执行级随机状态 \(\xi\) 确定随机化的R-FGSM更新。由于梯度从当前图像重新计算，这些随机化更新会影响后续的优化轨迹。

令 \(p_{k,m}(c)\) 为成功生成条件下归一化的输出分布，且多次调用间具有独立的新随机性，令 \(\mathcal{C} = \bigcup_{k,m} \mathrm{supp}(p_{k,m}) \subseteq C_{\mathrm{fp}}\)。对于共享 \(k, m\) 的独立成功输出，
\begin{equation}
\mathrm{Pr}[c_1 = c_2 \mid k, m] = \sum_{c \in \mathcal{C}} p_{k,m}(c)^2 \le \varepsilon_{\mathrm{prob}}. \tag{14}
\end{equation}
附录C通过区分具有相同和不同初始图像的执行来推导此一致界。它捕获了由两种随机源导致的输出离散度，并提供了下文使用的碰撞项。

\textbf{2) 多语义：}
对于生成的图像 \(c \in \mathcal{C}\) 和目标消息 \(m \in \mathcal{M}\)，定义兼容密钥集
\begin{equation}
\mathcal{K}(c, m) = \{ k' \in \mathcal{K} : \mathrm{DEC}_{k'}(c) = m \}. \tag{15}
\end{equation}
\textbf{命题（多语义）。}假设 \(C_{\mathrm{fp}} \subseteq \mathcal{K}\)。则对于每个 \(c \in \mathcal{C}\) 和 \(m \in \mathcal{M}\)，\(\mathcal{K}(c, m) \neq \emptyset\)。

\textbf{证明。}由于 \(c \in C_{\mathrm{fp}} \subseteq \mathcal{K}\)，将 \(c\) 用作图像密钥，并选择一个成功输出 \(c' \leftarrow \mathrm{ENC}_c(m)\)。正确性给出 \(F_c(c') = \mathrm{Et}(m)\)。由于融合图像加法是可交换的，
\begin{equation}
F_{c'}(c) = F_c(c') = \mathrm{Et}(m). \tag{16}
\end{equation}
由于 \(c' \in C_{\mathrm{fp}} \subseteq \mathcal{K}\)，\(\mathrm{DEC}_{c'}(c) = m\)，所以 \(c' \in \mathcal{K}(c, m)\)。\(\square\)

这种构造为每个选定消息提供了对单张图像的替代解释。兼容密钥特定于所选的 \((c, m)\) 对。下文的分析考虑了共享密钥的记录之间的依赖关系。

通过引入双重随机源并界定碰撞概率，结合融合映射的可交换性，
我们证明了同一消息多次生成满足特定的碰撞上界，且单张图像存在多语义解释。
这一结论说明，攻击者既无法通过收集多次通信记录来寻找重复模式，也无法在截获单张密文图像后确定其唯一对应的明文，从而为抵抗统计分析和唯密文攻击提供了理论支撑。
\subsubsection{形式化安全分析}

DNN、编码规则和生成设置是公开的。实例化配置固定了KeyGen、\(r\)和\(\xi\)的分布、完整的生成过程以及数值设置。每个游戏采样一个\(k \leftarrow \mathrm{KeyGen}()\)，并使用上述成功输出分布和新的生成随机性。查询计数在该密钥下跨消息和会话累计。对手是概率多项式时间算法，游戏涉及传输的图像和指定的预言机响应。

\textbf{1) IND-CPA优势界：} 对手\(\mathcal{A}\)在挑战前后自适应地查询\(\mathrm{ENC}_k\)。它提交等长消息\(m_0, m_1\)，对于均匀的\(b \in \{0,1\}\)，接收\(c^* \leftarrow \mathrm{ENC}_k(m_b)\)，并输出\(b'\)。其优势为\(\mathrm{Adv}_{\mathcal{A}}^{\mathrm{IND-CPA}} = |\Pr[b' = b] - 1/2|\)。令\(q\)统计所有生成输出，包括\(c^*\)。

我们将真实游戏与参考游戏进行比较，参考游戏为独立采样的均匀消息生成每个输出。令\(\eta_q\)为在最多\(q\)个输出且具有正参考似然的碰撞自由记录上，其转录似然比偏离1的最大绝对偏差。它衡量了这些记录中剩余的消息依赖性；附录D给出了其显式定义。

\textbf{定理1（查询有界IND-CPA优势）。}每个接收最多\(q\)个输出的此类对手满足
\begin{equation}
\mathrm{Adv}_{\mathcal{A}}^{\mathrm{IND-CPA}} \le \frac{\eta_q}{2} + \binom{q}{2} \varepsilon_{\mathrm{prob}}. \tag{17}
\end{equation}
混合证明利用正确性将真实生成表示为来自参考分布的条件采样。碰撞自由记录对统计距离的贡献至多为\(\eta_q/2\)，而重复输出对统计距离的贡献至多为\(\binom{q}{2}\varepsilon_{\mathrm{prob}}\)。参考游戏不泄露关于\(b\)的任何信息。

对于优势预算\(\varepsilon_{\mathrm{sec}} = \varepsilon_{\mathrm{mix}} + \varepsilon_{\mathrm{col}} < 1/2\)，附录D规定了关于\(\eta_Q\)和\(\varepsilon_{\mathrm{prob}}\)的条件，以确定可接受的上限\(Q \ge 1\)。对于满足这些条件的任何上限，
\begin{equation}
\mathrm{Adv}_{\mathcal{A}}^{\mathrm{IND-CPA}} \le \varepsilon_{\mathrm{sec}}, \quad 1 \le q \le Q. \tag{18}
\end{equation}
这给出了CCA1分析的仅生成基线。可接受的上限取决于实例化分布和所选择的优势预算。

\textbf{2) IND-CCA1 优势界：} 对手 \(\mathcal{A}'\) 可以在挑战前自适应地查询 \(\mathrm{ENC}_k\) 和 \(\mathrm{DEC}_k\)，但挑战后只能查询 \(\mathrm{ENC}_k\)。挑战消息与优势定义沿用上述 IND-CPA 游戏。保留 \(q\) 作为生成输出的数量，并令 \(q_d\) 统计挑战前的恢复查询 \(x \in C_{\mathrm{fp}}\)。

为了分离恢复访问的影响，比较同密钥和分裂密钥预言机游戏。两者都使用 \(k\) 进行生成；分裂密钥游戏在一个独立的固定密钥 \(\bar{k}\) 下回答首次非重放恢复查询，同时保留缓存和精确重放响应。令 \(\delta_{\mathrm{dec}}(q, q_d)\) 为这些限制下它们的最大区分概率差距；见附录 E。

\textbf{定理 2（查询有界 IND-CCA1 优势）。} 对于每个这样的 \(\mathcal{A}'\)，存在一个 IND-CPA 对手 \(\mathcal{A}\) 使得
\begin{equation}
\begin{aligned}
\mathrm{Adv}_{\mathcal{A}'}^{\mathrm{IND-CCA1}} &\le \mathrm{Adv}_{\mathcal{A}}^{\mathrm{IND-CPA}} + \delta_{\mathrm{dec}}(q, q_d) \\
&\le \frac{\eta_q}{2} + \binom{q}{2} \varepsilon_{\mathrm{prob}} + \delta_{\mathrm{dec}}(q, q_d).
\end{aligned} \tag{19}
\end{equation}

分裂密钥视图可以使用 CPA 预言机和 \(k\) 下的本地恢复来模拟。附录 E 通过密钥解耦行为以及较早的恢复输入后来作为生成输出出现的概率，来界定额外的差距。给定一个可接受的 \(Q\)，它指定了恢复查询限制 \(q_d \ge 0\) 的条件，使得
\begin{equation}
\mathrm{Adv}_{\mathcal{A}'}^{\mathrm{IND-CCA1}} \le \varepsilon_{\mathrm{sec}} + \varepsilon_{\mathrm{dec}}, \quad 1 \le q \le Q, \quad q_d \le Q_d, \tag{20}
\end{equation}
其中 \(\varepsilon_{\mathrm{sec}} + \varepsilon_{\mathrm{dec}} < 1/2\)。当 \(q_d = 0\) 时，额外差距消失，结果退化为定理 1。

这些条件规定了允许的每密钥使用量：部署可以强制执行不大于可接受 \((Q, Q_d)\) 的生成和恢复上限，同时保持所分析的预言机访问模式。有用限制的存在性和大小取决于实例化分布。在部署中还应通过密钥更新和访问限制控制累计查询。

我们严格推导出了查询有界的 IND-CPA 与 IND-CCA1 优势上界，并确定了安全成立的条件。这意味着，在满足相关输出分布条件，且遵守给定的每密钥生成查询上限($Q$)与恢复查询上限($Q_d$)
时，攻击者的区分优势被严格限制在预定的安全预算之内。该结论表明系统可以通过限制访问频次和定期更新密钥，来确保在长期运行中始终维持给定的安全性标准。